% Round 18. ci_global branch. SG18 candidate; independent review required.
\section{An explicit global surface ambient and complete height families}
\label{r18sg:section}

The supplied global surface hypothesis of SC18 is discharged here.
The construction has one fixed number field, a dense full section
space, and a deep degeneration at the fixed rational prime \(5\).
The local surface at that prime has degree-two cohomology which
is not generated by geometric divisors, even after any finite
residue-field extension.  Thus the new SC18 curve-comparison
interface is used essentially; the sufficient all-algebraic
\(H^2\) condition in PG18 is not substituted for it.

The arithmetic input about complex multiplication and good
reduction is the original paper of Serre--Tate:
\emph{Good reduction of abelian varieties}, Annals of Mathematics
\textbf{88} (1968), 492--517, Theorem 1 on printed page 493 and
Theorem 6(a) on printed page 503.  These two statements were
checked in the primary source.  In the scanned PDF, including
its cover, they are respectively pages 3 and 13.
\begin{quote}
\url{https://wstein.org/papers/bib/Serre-Tate-Good_Reduction_of_Abelian_Varieties.pdf}
\end{quote}
We use only their precise implications: CM gives potential good
reduction at finite places; over a discrete valuation field with
perfect residue field, trivial prime-to-residue inertia on a
Tate module gives good reduction.  The finite congruence-subgroup
argument and the field arrangements below are written out.

\subsection{One explicit number field with an everywhere good elliptic curve}

Put
\[
 E_0/\mathbf Q:\quad y^2=x^3-x,\qquad
 F_{\rm tor}=\mathbf Q(E_0[3]).
 \tag{SG18.1}\label{r18sg:torsion-field}
\]
The notation \(F_{\rm tor}\) refers to the field generated by
the coordinates of every geometric \(3\)-torsion point, rather
than the field of one chosen point.

\begin{lemma}[The good surface input, with control at \(5\)]
\label{r18sg:global-good}
The finite Galois field \(F_{\rm tor}/\mathbf Q\) is unramified
at \(5\).  The elliptic curve \(E_0\) has good reduction at every
finite place of \(F_{\rm tor}\), and extends to a smooth
projective elliptic scheme
\(\mathcal E/\mathcal O_{F_{\rm tor}}\).
Consequently
\[
 \mathcal S=\mathcal E\times_{\mathcal O_{F_{\rm tor}}}\mathcal E,
 \qquad
 \mathcal A=
   \mathcal O_{\mathcal E}(3[0])
      \boxtimes\mathcal O_{\mathcal E}(3[0])
 \tag{SG18.2}\label{r18sg:surface}
\]
is a smooth projective surface model with geometrically
integral fibers at every finite place, and \(\mathcal A\)
is relatively very ample with fiberwise self-intersection \(18\).
\end{lemma}

\begin{proof}
The displayed Weierstrass equation has discriminant \(64\).
It therefore gives good reduction at every odd prime.
Over \(\mathbf Q(i)\), the map
\((x,y)\mapsto(-x,iy)\) is an automorphism whose square is
the elliptic inverse.  It gives complex multiplication
by \(\mathbf Q(i)\).  Serre--Tate, Theorem 6(a), gives
potential good reduction at the places over \(2\), hence
also potential good reduction of \(E_0/\mathbf Q_2\).

Write \(T_3(E_0)\simeq\mathbf Z_3^2\).
Potential good reduction implies that the image of inertia
at \(2\) in \(\operatorname{GL}_2(\mathbf Z_3)\) is finite:
over a finite extension with good reduction its inertia
subgroup acts trivially by Serre--Tate, Theorem 1, and this
subgroup has finite index in the original inertia group.
After passage to any completion of \(F_{\rm tor}\) over
\(2\), inertia also acts trivially modulo \(3\).  Its image
is therefore a finite subgroup of
\[
                    1+3M_2(\mathbf Z_3).
 \tag{SG18.3}\label{r18sg:congruence-group}
\]

For clarity, this group has no nonidentity element of finite
order.  If \(g=1+3^aB\ne1\), where \(a\ge1\) and some entry
of \(B\) is a \(3\)-adic unit, and if \(q\ne3\) is prime,
the first term of
\[
                  g^q-1=q3^aB+\sum_{j=2}^q
                            \binom qj3^{aj}B^j
\]
has an entry of valuation \(a\), whereas every subsequent
term has all entries of valuation at least \(2a\).
Thus \(g^q\ne1\).  For \(q=3\), the first term has an entry
of valuation \(a+1\); the remaining two terms have entrywise
valuations at least \(2a+1\) and \(3a\), both strictly
greater than \(a+1\).  Again \(g^3\ne1\).
A nontrivial finite-order element would have a power of
prime order, a contradiction.  The finite inertia image
in \eqref{r18sg:congruence-group} is trivial.
Serre--Tate, Theorem 1, now proves good reduction at every
place over \(2\).  At all other places good reduction of
the original curve persists under base change, including
the places over \(3\).

At \(5\), the original smooth elliptic model has finite
etale \(3\)-torsion.  Inertia acts trivially on all of
\(E_0[3]\).  Since the full torsion field is the fixed
field of the kernel of this finite representation,
\(F_{\rm tor}/\mathbf Q\) is unramified at \(5\).

We also record why the output is the projective model
required by SC18.  Over each local ring of
\(\mathcal O_{F_{\rm tor}}\) the good elliptic model is the
proper Neron model.  These local models are uniquely
isomorphic on their overlaps with the generic elliptic
curve: a smooth proper group model has the Neron
extension property.  Equivalently, take the global
Neron model and use its properness at every local ring.
Finite presentation lets the finitely many models
different from the original model over
\(\mathcal O_{F_{\rm tor}}[1/2]\) spread to neighborhoods.
Uniqueness patches them to a proper smooth elliptic
scheme \(\mathcal E\) over the whole number ring.
This uses the standard Neron-model interpretation of
good reduction, also used in the proof of the cited
Serre--Tate theorem; it does not assert that one global
Weierstrass equation has unit discriminant.

The identity section is an effective Cartier divisor.
On every geometric genus-one fiber, a line bundle of
degree \(3\) separates every length-two subscheme:
by Serre duality the relevant \(H^1\) is dual to
sections of a line bundle of degree \(-1\), and hence
vanishes.  Its own \(H^1\) vanishes as well.
Cohomology and base change and the evaluation map
therefore give a relative closed immersion by
\(\mathcal O_{\mathcal E}(3[0])\).
The elliptic scheme is projective, with geometrically
integral fibers.  The external product of these
embeddings proves the corresponding assertions for
\(\mathcal S\) and \(\mathcal A\).
Finally
\(\deg_{\mathcal S}(\mathcal A^2)=2\cdot3\cdot3=18\)
on every geometric fiber.
\end{proof}

\subsection{Actual coefficients and the full section space}

For the remainder put
\[
 \mathcal Y=\mathcal S\times\mathbf P^3,\qquad
 \mathcal L=\mathcal A\boxtimes\mathcal O_{\mathbf P^3}(1).
 \tag{SG18.4}\label{r18sg:ambient}
\]
All products in this subsection are over
\(\mathcal O_{F_{\rm tor}}\).
No freeness of a Hodge or section module over this
Dedekind ring is assumed.

\begin{lemma}[Free ambient coefficients and projective normality]
\label{r18sg:coefficients}
There is a closed immersion
\(\mathcal Y\hookrightarrow\mathbf P^N_{\mathcal O_{F_{\rm tor}}}\)
with \(\mathcal L=\mathcal O(1)|_{\mathcal Y}\).
The restrictions of degree-\(d\) homogeneous polynomials
in its \(N+1\) coordinates span \(H^0(Y_k,\mathcal L^d)\)
for every field fiber \(k\) and every \(d\ge1\).
The same restriction map onto the relative section
module is surjective over the number ring.
\end{lemma}

\begin{proof}
The section module of \(\mathcal O_{\mathcal E}(3[0])\)
is locally free of rank \(3\), by the genus-one
vanishing used above.  Choose finitely many generators
of this module over the number ring.  A finite free
module surjecting onto it embeds its projective bundle
as a closed subscheme of an ordinary relative
projective space.  Thus this choice provides an
ordinary relative projective embedding of
\(\mathcal E\), without trivializing its rank-three
section bundle.  Do this for the two elliptic
factors, and combine with the standard embedding
of \(\mathbf P^3\) by the Segre map.

Over any field, the degree-three elliptic embedding
is a plane cubic, and the exact sequence
\[
 0\longrightarrow\mathcal O_{\mathbf P^2}(d-3)
  \longrightarrow\mathcal O_{\mathbf P^2}(d)
  \longrightarrow\mathcal O_E(3d[0])
  \longrightarrow0
\]
shows that degree-\(d\) forms span all sections:
\(H^1(\mathbf P^2,\mathcal O(t))=0\) for every integer
\(t\).  The extra free generators merely linearly
embed this plane in a larger projective space.
The two elliptic factors and the projective factor
therefore have this spanning property in every degree.
Every tensor of three degree-\(d\) monomials can be
paired factor by factor into a product of \(d\)
Segre coordinates.  It follows that the Segre
degree-\(d\) forms span the entire section space
on \(Y_k\).

For \(d\ge1\), each positive-degree line bundle on
an elliptic fiber has no \(H^1\), and the positive
twists on \(\mathbf P^3\) have no higher cohomology.
Kunneth and cohomology and base change show that
the relative section module for \(\mathcal L^d\)
is locally free and commutes with field base change.
The cokernel of its polynomial restriction map
is a finite module whose fiber at every maximal
ideal is zero by the preceding calculation.
Nakayama's lemma gives the asserted global
surjectivity.  All residue congruences below
can consequently be imposed in finite free
modules of actual ambient polynomials.
\end{proof}

\subsection{The full-height family with a fixed deep prime}

\begin{theorem}[An explicit surface-ambient complete-height family]
\label{r18sg:height-family}
There is a fixed finite extension \(F/F_{\rm tor}\),
unramified over \(F_{\rm tor}\) at every place over \(5\),
and a Zariski-dense family of smooth geometrically
integral divisors
\[
                     X\in|\mathcal L^5|_{Y_F}
 \tag{SG18.5}\label{r18sg:linear-system}
\]
with the following properties.
\begin{enumerate}
\item For every finite extension \(F'/F\), every finite
      place \(w\) of \(F'\), every codimension, and every
      specified geometrically homologically trivial
      rational cycle on \(X_{F'_w}\), strong G holds
      with the exact specified support and all compact
      ordinary actual tests, for the given absolute value.
\item The entire complementary homologically trivial
      Chow groups of \(X_{F'}\) have the BB pairing of
      PH16, for every finite extension \(F'/F\).
      In dimension four the complementary degrees are
      \(r\) and \(5-r\), \(1\le r\le4\).
      GA17 supplies the regular-model admissible lifts
      on these full groups.
\item At every place of \(F\) over \(5\), the family has
      an actual five-factor product smoothing with
      \(\operatorname{rank}N^4=179\).
      Moreover \(b_1(X)=4\) and \(b_2(X)=7\).
      Hence no member is geometrically isomorphic to
      a smooth projective-space complete-intersection
      fourfold.  It is not potentially of good reduction
      at these places, and every strict semistable
      dominant witness there has crossing depth at
      least four.
\end{enumerate}
The density is in the full section space over this one
fixed field \(F\).  The theorem does not assert existence
over \(\mathbf Q\), or G for every member of the linear
system.
\end{theorem}

\begin{proof}
\emph{Step 1: reference and finite residue fields.}
Choose an integral ambient degree-five polynomial
\(H_0\) whose restriction to \(Y_{F_{\rm tor}}\)
has a smooth zero divisor.  Such polynomials exist
by the first-jet incidence proof and density of
an integral coefficient lattice in characteristic
zero.  Let \(S_0\) be the finite set of rational
primes different from \(5\) below the places where
this divisor is not smooth, together with \(2\).
At each place \(v\) of \(F_{\rm tor}\) over every
\(p\in S_0\), choose a smooth degree-five section
over some finite extension \(\kappa(v)'\) of
\(\kappa(v)\).  This is done at \emph{all} such
places, including any where \(H_0\) was already
smooth, since a rational coefficient divisible
by \(p\) vanishes at every one of them.
The smooth-section open is geometrically nonempty,
so it has a point over a finite extension of the
finite residue field; no fixed-degree finite-field
Bertini assertion over the original field is used.

At every place \(v\mid5\), instead choose a finite
residue extension and five sections \(f_{1,v},\ldots,
f_{5,v}\) of \(\mathcal L\), together with a
section \(Q_v\) of \(\mathcal L^5\), satisfying all
the geometric transversality hypotheses of PG18.
The required simultaneous parameter open is
geometrically nonempty by Bertini on the smooth
fivefold \(Y_{\overline{\kappa(v)}}\): the factors
have geometric simple normal crossings, all
expected positive-dimensional intersections are
smooth, and \(Q_v\) is transverse to all of them
and avoids the finitely many top points.
Every positive-dimensional smooth ample complete
intersection is connected, by degree-zero weak
Lefschetz applied successively, and therefore is
geometrically integral.  This also checks the
positive-dimensional intersections involving
\(Q_v\).  The finite-field point of the joint open
exists after some finite residue extension.

Take \(N\) divisible by the degrees of all the
chosen finite residue extensions over their
original \(\kappa(v)\).  At the finitely many
selected maximal ideals of \(\mathcal O_{F_{\rm tor}}\),
choose monic irreducible polynomials of degree \(N\).
Coefficientwise CRT gives one monic polynomial
\(g\in\mathcal O_{F_{\rm tor}}[T]\) with all these
reductions.  It is irreducible over \(F_{\rm tor}\):
any monic factor over its fraction field would
have coefficients integral over the integrally
closed ring \(\mathcal O_{F_{\rm tor}}\), and
therefore in that ring, contradicting one of
the irreducible reductions.  Put
\[
                 F=F_{\rm tor}[T]/(g).
\]
Its localization at each selected place is
unramified and inert of degree \(N\).
Indeed, the reduction of \(g\) is irreducible
and separable, so its local polynomial algebra
is finite etale with the indicated residue
field and already is the corresponding
integral closure.  Every selected residue
extension is now contained in the residue
field of \(F\).  In particular \(F/\mathbf Q\)
is unramified at \(5\), since \(F_{\rm tor}/
\mathbf Q\) already was.

\emph{Step 2: integral sections and a Lefschetz pencil.}
Over \(\mathcal O_F\), lift the chosen factor and
smoothing polynomials at all places over \(5\)
simultaneously by CRT in the free coefficient
modules of Lemma~\ref{r18sg:coefficients}.
Write their lifts \(f_1,\ldots,f_5,Q\) and
\(P=\prod_i f_i\).
The same coefficient lattices meet the generic
joint Bertini open, if it is imposed as well.
Lift the chosen smooth sections at the other
selected places to residue polynomials \(G_v\).
Put
\[
 B=\prod_{p\in S_0}p,\qquad M=25B.
 \tag{SG18.6}\label{r18sg:M}
\]
Choose an integral ambient polynomial \(H_1\)
subject to
\[
 H_1\equiv P+5Q\pmod{25},\qquad
 H_1\equiv G_v\pmod{\mathfrak p_v}
       \quad(v\mid S_0),
 \tag{SG18.7}\label{r18sg:congruences}
\]
and require that the pencil \([uH_0+vH_1]\)
belong to the fixed-\(H_0\) Lefschetz open.
These congruences are pairwise compatible:
the places over \(5\) are disjoint from those
over \(S_0\), and \(25\mathcal O_F\) is the
product of the squared ideals over \(5\)
because \(F/\mathbf Q\) is unramified there.
Their solution set is a nonempty translate
of a full integral coefficient lattice,
hence is Zariski dense over \(F\).

The fixed-member open is nonempty on this
actual ambient, as proved in CG18 and used
with this scope in SC18.  Its proof needs
only one-point first and second jets and
independent first jets at two distinct
points, supplied by degree-five polynomials
in any projective embedding.  On a smooth
fivefold these give a discriminant divisor
whose general point is a unique nondegenerate
critical point.  Projection of the
discriminant from the fixed smooth section
is finite and generically etale in
characteristic zero.  A general line
through that section has only transverse
one-node singular members and a smooth
base locus.  Thus the congruence lattice
indeed contains the required \(H_1\);
no projective-space-only hypersurface
theorem is applied to \(Y\).

After fixing \(H_1\), spread out this pencil
and enlarge a finite exceptional set so
that its fibers have good or one-node
reduction at every other place.  Include
all places above \(S_0\cup\{5\}\) in the
exceptional set, and let \(T\) be the
finite set of underlying rational primes.
For each integer \(a\) satisfying
\[
 a\equiv1\pmod{25},\qquad a\equiv1\pmod B,
 \qquad
 a\equiv0\pmod p\quad
      (p\in T\setminus(S_0\cup\{5\})),
 \tag{SG18.8}\label{r18sg:a}
\]
let \(X_a\) be the divisor
\[
                         MH_0+aH_1=0
                     \quad\hbox{on }Y_F .
 \tag{SG18.9}\label{r18sg:family}
\]
The finitely many singular generic members
of the pencil are omitted.  Every remaining
member is smooth and geometrically integral.

\emph{Step 3: actual models at every finite place.}
At every place over \(5\), write
\(H_1=P+5Q+25R\).  Since \(a\) is a unit,
divide \eqref{r18sg:family} by \(a\) to obtain
the exact local equation
\[
 P+5Q_a=0,\qquad
 Q_a=Q+5R+5B a^{-1}H_0,\qquad
                    \overline{Q_a}=\overline Q .
 \tag{SG18.10}\label{r18sg:exact-product}
\]
Here \(5\) is the native uniformizer.
Thus this is literally a PG18 smoothing
with the previously chosen transverse
special factors and smoothing section.
At the places over \(S_0\), \(M\) vanishes
and \(a\equiv1\); the reduction is the
chosen smooth section \(G_v\).
At any remaining exceptional place,
\(M\) is a unit and \(a\equiv0\);
the reduction is a nonzero multiple of
the smooth reference \(H_0\).
Every other place lies in the spread
Lefschetz locus and has a smooth or
ordinary quadratic special divisor.
All characteristic-two places were
already forced to be smooth.

For completeness, these are the actual
flat integral projective models needed
by the local theorems.  The ambient
\(\mathcal Y\) is smooth with integral
geometric fibers, and each displayed
special defining section is nonzero.
It is therefore a nonzerodivisor on
the special ambient.  The relative
Cartier divisor is flat over the DVR:
if \(\pi x\) lies in the defining ideal,
reduction modulo \(\pi\) and cancellation
of the nonzerodivisor shows that \(x\)
does as well.  Its smooth integral
generic fiber and this torsion-freeness
make the total divisor integral.

The smooth and isolated-node special
divisors are geometrically integral.
Indeed,
\[
 H^1(S_{\bar k}\times\mathbf P^3_{\bar k},
                   A^{-5}\boxtimes\mathcal O(-5))=0,
\]
because the two Kunneth terms contain
\(H^0(\mathbf P^3,\mathcal O(-5))=0\)
or \(H^1(\mathbf P^3,\mathcal O(-5))=0\).
The divisor sequence gives
\(H^0(X_{\bar k},\mathcal O)=\bar k\).
An isolated-node Cartier fourfold is
\(S_2\) and \(R_1\), hence normal;
connected normal components cannot
be distinct.  The same argument
applies to smooth fibers.

\emph{Step 4: all cycles, finite extensions, and heights.}
At every place over \(5\), the product
surface \(S\) gives the uniform curve
comparison of CD17, independently
of its divisor rank.  SC18 propagates
it through the two ample cuts to
the initial double threefolds and
then through the actual codimension-two
edge blowups.  PG18 gives \(U_2=0\)
on this model and, with that curve
comparison, full local G in dimension
four.  This reasoning reapplies
directly to every new cycle over
every finite unramified extension;
RI18 gives the conclusion for every
finite valued extension, including
ramified ones.  The exact cycle
support, all actual tests, and the
original absolute value are retained
by those frozen interfaces and RI18.
No new cycle is assumed to descend
from a smaller field.

At every other finite place the
good-reduction theorem or ON17
applies to the literal model just
checked, with its extension statement
for all new cycles.  Therefore every
finite place of every finite \(F'/F\)
has the complete local G domain.
The global hypotheses of PH16 are
now all satisfied in each complementary
degree; it supplies the claimed BB
pairings.  GA17 supplies admissible
Chow lifts on precisely this proved
full height domain.

\emph{Step 5: density over a fixed field.}
All residue-point and field choices
were made before varying \(H_1\).
The resulting fixed congruence lattice
of allowed \(H_1\)'s is Zariski dense
in the ambient coefficient space,
and its intersection with the
fixed-member open remains dense.
For each such \(H_1\), there are
infinitely many integers \(a\) in
the progression \eqref{r18sg:a}
outside the finite set of singular
generic pencil parameters.
A polynomial vanishing on every
constructed coefficient vector
\(MH_0+aH_1\) therefore vanishes
on the whole affine line with
that direction.  In particular
it vanishes at \(MH_0+H_1\).
Varying \(H_1\) in its dense
lattice proves that the polynomial
is zero.  Lemma~\ref{r18sg:coefficients}
surjects this coefficient space
onto the full section space,
so the asserted projective density
in \eqref{r18sg:linear-system}
follows.

\emph{Step 6: the distinguishing invariants.}
Put \(h=c_1(\mathcal L)\) and let
\(b\) denote the projective
hyperplane on \(\mathbf P^3\).
Since \(h=A+b\),
\[
             \deg_Y h^5=\binom52\deg_S A^2=180 .
 \tag{SG18.11}\label{r18sg:top-degree}
\]
The five factor intersection has
180 geometric points.  All lower
positive-dimensional strata are
connected.  The HM18/NG18 extreme
weight calculation, applied to the
actual ordered product resolution,
therefore gives
\[
                      \operatorname{rank}N^4=180-1=179 .
 \tag{SG18.12}\label{r18sg:monodromy}
\]
The calculation retains the native
closed points and their full degrees;
180 is their total geometric count.
Weak Lefschetz gives
\(H^j(X_{\bar F},\mathbf Q_\ell)
 =H^j(Y_{\bar F},\mathbf Q_\ell)\)
for \(j<4\).  Two elliptic factors
give \(b_1(S)=4\) and \(b_2(S)=6\);
the projective factor adds one to
the second Betti number.  Thus
\(b_1(X)=4\) and \(b_2(X)=7\).
A smooth projective-space complete
intersection of dimension four has
\(b_1=0,b_2=1\), again by weak
Lefschetz.  These are intrinsic
geometric invariants, not statements
only about the chosen polarization.
Finally the nonzero fourth monodromy
persists under finite extension and
gives the asserted potential-good
and crossing-depth exclusions by
NM17.
\end{proof}

\subsection{Why the curve gate adds a real new input}

At \(5\), the original equation gives
\(\#E_0(\mathbf F_5)=8\), so its
Frobenius trace is \(-2\) and the
two degree-one eigenvalues are
\(\alpha=-1+2i\), \(\beta=-1-2i\).
The square surface has the two
Tate-twisted degree-two eigenvalues
\[
                       \alpha^2/5=(-3-4i)/5,\qquad
                       \beta^2/5=(-3+4i)/5 .
 \tag{SG18.13}\label{r18sg:nonalgebraic}
\]
Neither is a root of unity:
the roots of unity in \(\mathbf Q(i)\)
are \(1,-1,i,-i\).  A geometric
divisor over a finite field is
defined over a finite extension,
so its Tate-twisted class is
fixed by some positive Frobenius
power.  The two displayed
eigenspaces are consequently
outside the span of divisor
classes.  Every residue extension
in the preceding construction
raises the displayed eigenvalues
to a positive power, which
cannot turn a non-root of unity
into a root of unity.
The projective factor adds only
an algebraic line in \(H^2\).

Thus at the very prime where the
family has deep degeneration,
the ambient's \(H^2(1)\) is
not all algebraic.  Its curve
comparison nevertheless holds,
because the projective-bundle
formula over a surface needs
only divisor numerical-to-homological
comparison and degree-zero
zero-cycle comparison, and SC18
transfers exactly that property
to the double strata.
No Tate conjecture for arbitrary
codimension was invoked.

\paragraph{Exact contribution and limits.}
The constructed fourfolds have
full homological Chow domains,
not merely a pullback or product
cycle subspace, at every finite
place and after every finite
number-field extension.
They have intrinsic \(b_2=7\),
so they differ from both the
complete-intersection families
of CG18/CE18 and the \(b_2=2\)
product-ambient family of PA18.
This distinction makes no claim
about disjointness from every
earlier motivic sufficient sector.
In particular an ample divisor
of an abelian-surface product
is not asserted to have an
abelian motive merely from its
ambient.
The existence and density here
do not imply G for an arbitrary
variety or an arbitrary member
of this linear system.
Those general assertions and
the corresponding arbitrary-variety
complete unconditional height
problem remain open.
